% Public-source biographies authorized by the user on 31 October 2026.
% Evidence: docs/submission_successor_20261031/AUTHOR_PUBLIC_SOURCES.json.
% Keep author order. No photographs or corresponding-author designation.

% J2/J3: self-entered ORCID education and Uniandes thesis; no degree-conferral
% date, master's degree, or current university employment is asserted.
\begin{IEEEbiographynophoto}{Jorge Luis Mayorga Taborda}
studied electronic engineering at Universidad de los Andes, Bogot\'a, Colombia. His academic work has addressed distributed control using population dynamics, including applications to urban traffic systems.\footnote{Biographical drafting assistance: OpenAI Codex~\cite{openaiCodex2026}.}
\end{IEEEbiographynophoto}

% Y3/Y4: Uniandes thesis and IEEE-deposited publication metadata.
% No degree-conferral claim, current job title, or middle-name addition.
\begin{IEEEbiographynophoto}{Yessica Africano Rodriguez}
carried out undergraduate thesis work on photovoltaic integration at Universidad de los Andes, Bogot\'a, Colombia. Her published research includes co-simulation and stochastic assessment of photovoltaic hosting capacity in distribution systems.
\end{IEEEbiographynophoto}

% M1/M3/M4: Pitt institutional profile. Rank intentionally omitted because
% its contact card says Associate Professor while its narrative says assistant.
\begin{IEEEbiographynophoto}{Masoud Barati}
received the Ph.D. degree in electrical engineering from the Illinois Institute of Technology in 2013. He is with the Department of Electrical and Computer Engineering, University of Pittsburgh, Pittsburgh, PA, USA. His research includes mathematical modeling and optimization of electric power systems.
\end{IEEEbiographynophoto}

% P1/P2/P3: current GWU faculty page.
\begin{IEEEbiographynophoto}{Payman Dehghanian}
received the Ph.D. degree from Texas A\&M University in 2017. He is an Associate Professor with the Department of Electrical and Computer Engineering, The George Washington University, Washington, DC, USA. His research includes power-system reliability, resilience, and data-informed decision-making.
\end{IEEEbiographynophoto}

% D1/D2/D3: La Sabana institutional profiles; broad institutional degree wording.
\begin{IEEEbiographynophoto}{David Felipe Celeita Rodr\'iguez}
received the Ph.D. degree in engineering from Universidad de los Andes, Bogot\'a, Colombia. He is an Associate Professor and Director of Outreach and Engagement with the Faculty of Engineering, Universidad de La Sabana, Ch\'ia, Colombia. His research includes power-system protection, fault location, and real-time simulation.
\end{IEEEbiographynophoto}
